\title{DeepSeekMath: Pushing the Limits of Mathematical Reasoning in Open Language Models}

\begin{abstract}
The structured and intricate nature of mathematical reasoning presents considerable difficulties for language models. In this study, we present DeepSeekMath~7B, which is built by further pre-training DeepSeek-Coder-Base-v1.5 7B using 120B math-focused tokens collected from Common Crawl, as well as natural language and code datasets. DeepSeekMath~7B attains a notable 51.7\% on the MATH benchmark, which comprises competition-level problems, all achieved without the support of external tools or voting mechanisms—bringing it close to the performance of Gemini-Ultra and GPT-4. When evaluating self-consistency across 64 sampled responses, DeepSeekMath~7B achieves a score of 60.9\% on MATH. This advancement in mathematical reasoning is primarily attributed to two aspects: Firstly, we effectively utilize publicly available web content via a carefully crafted data selection workflow. Secondly, we propose Group Relative Policy Optimization (GRPO)—an adaptation of Proximal Policy Optimization (PPO)—that not only strengthens mathematical reasoning but also yields improved memory efficiency compared to standard PPO.
\end{abstract}

\section{Introduction}

Large language models (LLMs) have transformed how artificial intelligence addresses mathematical reasoning tasks, leading to notable progress on benchmarks for quantitative reasoning \citep{MATH} and geometry reasoning \citep{trinh2024solving}. In addition, these systems have demonstrated significant utility in helping individuals tackle challenging mathematical problems \citep{tao}. Nevertheless, the most advanced models, including GPT-4 \citep{gpt4} and Gemini-Ultra \citep{gemini}, remain inaccessible to the public, while available open-source alternatives exhibit a substantial gap in performance.

This paper presents DeepSeekMath, a specialized language model in the mathematics domain that notably exceeds the performance of existing open-source systems and narrows the gap with GPT-4 on academic benchmarks. Central to this advance is the construction of the DeepSeekMath Corpus, a large-scale, high-quality pre-training dataset comprising 120B mathematical tokens. The dataset is assembled from Common Crawl (CC) resources through a classifier based on fastText \citep{joulin2016fasttext}. Initially, the classifier is trained with positive samples drawn from OpenWebMath \citep{openwebmath} and a wide assortment of unrelated web content as negative samples. The trained classifier is then applied to identify further positive examples within CC data, which are subsequently verified and refined by human annotators. These curated samples are used to retrain the classifier, thereby enhancing its accuracy. Evaluations demonstrate that the resulting corpus sustains robust performance: our foundational model, DeepSeekMath-Base 7B, secures 64.2\% accuracy on GSM8K \citep{gsm8k} and 36.2\% on the MATH competition dataset \citep{MATH}, surpassing Minerva 540B \citep{minerva}. Additionally, the multilingual breadth of the DeepSeekMath Corpus yields notable improvements in Chinese mathematics benchmarks \citep{wei2023cmath,agieval}. We regard our methodology for mathematical data curation as a constructive foundation for future community-driven research, acknowledging considerable potential for further enhancement.

DeepSeekMath-Base is built upon DeepSeek-Coder-Base-v1.5 7B \citep{deepseek-coder}, as initializing from a code-oriented model proves more advantageous than beginning with a standard language model. Additionally, our findings reveal that mathematical training bolsters the model’s performance not only in mathematical domains, but also on tasks from the MMLU \citep{mmlu} and BBH \citep{bbh} benchmarks. This suggests that such training strengthens both mathematical proficiency and broader reasoning abilities.

Following the pre-training phase, we perform mathematical instruction tuning on DeepSeekMath-Base utilizing datasets that incorporate chain-of-thought \citep{cot}, program-of-thought \citep{pot,pal}, and tool-integrated reasoning \citep{tora} approaches. The tuned model, DeepSeekMath-Instruct 7B, surpasses all other 7B models and achieves performance on par with instruction-tuned open-source models of 70B scale.

Additionally, we present Group Relative Policy Optimization (GRPO), a novel reinforcement learning (RL) approach derived from Proximal Policy Optimization (PPO) \citep{schulman2017proximal}. Distinct from PPO, GRPO eliminates the need for a critic model by computing the baseline via group scores, thereby markedly decreasing computational demands during training. Utilizing only a portion of English instruction tuning data, GRPO achieves notable gains over the already competitive DeepSeekMath-Instruct, delivering marked improvements for both in-domain benchmarks (GSM8K: 82.9\% $\rightarrow$ 88.2\%, MATH: 46.8\% $\rightarrow$ 51.7\%) and out-of-domain tasks (e.g., CMATH: 84.6\% $\rightarrow$ 88.8\%) in the reinforcement learning stage. We further provide a generalized framework for interpreting various techniques, such as Rejection Sampling Fine-Tuning (RFT) \citep{yuan2023scaling}, Direct Preference Optimization (DPO) \citep{dpo}, PPO, and GRPO, revealing through this lens that these approaches can all be understood as forms of direct or streamlined RL methods. We perform comprehensive experimental comparisons, including online versus offline training, output versus process-based supervision, and single-step versus iterative RL, to scrutinize critical aspects within this framework. Finally, we discuss the underlying reasons for performance gains in instruction-tuned models with RL and outline future prospects for developing more effective RL techniques in accordance with this unified conceptualization.

\subsection{Contributions}
We present scalable math pre-training techniques and investigate the application and assessment of reinforcement learning methods.

\textbf{Math Pre-Training at Scale}

\begin{itemize}[topsep=0pt]
    \item This study demonstrates that the widely available Common Crawl resource holds substantial mathematical data. Through a carefully engineered data filtering process, we have assembled the DeepSeekMath Corpus, comprising 120B tokens sourced from web pages selected for their mathematical relevance. This dataset is nearly seven times larger than the math-specific web pages utilized by Minerva \citep{minerva}, and nine times greater in size compared to the newly introduced OpenWebMath \citep{openwebmath}.

\item The DeepSeekMath-Base 7B, our foundational pre-trained model, delivers results on par with Minerva 540B \citep{minerva}, demonstrating that parameter count alone does not solely dictate mathematical reasoning abilities. Robust performance can also be attained by a more compact model trained on high-caliber data.

\item We present the outcomes of our math training experiments. Pretraining models on code before subjecting them to math-specific training enhances their performance on mathematical tasks, regardless of whether external tools are utilized. These results contribute to resolving the enduring question: \textit{does code training improve reasoning abilities?} Our results suggest that, specifically in the context of mathematical reasoning, it indeed does.

\item While utilizing arXiv papers for training is a widespread practice, particularly within mathematics-focused research, our results indicate that this approach yields no significant enhancement across any of the mathematical benchmarks evaluated in this study.
\end{itemize}

\textbf{Exploration and Analysis of Reinforcement Learning}

\begin{itemize}[topsep=0pt]
    \item We propose Group Relative Policy Optimization (GRPO), a reinforcement learning technique that operates without a critic model by using group scores to estimate the baseline. This approach yields substantial savings in training resources in comparison to Proximal Policy Optimization (PPO).
    \item We show that applying GRPO leads to marked improvements in the performance of our instruction-tuned system, DeepSeekMath-Instruct, utilizing only instruction-tuning datasets. Additionally, we report improved generalization to out-of-domain data observed throughout the reinforcement learning phase.
    \item We establish a comprehensive framework that unifies various strategies such as RFT, DPO, PPO, and GRPO. Within this framework, we carry out thorough experimental investigations, including comparisons between online and offline learning, outcome versus process-level supervision, as well as single-turn versus iterative reinforcement learning, among others, to elucidate key aspects of these approaches.
    \item Drawing from our unified framework, we examine the underlying factors that contribute to the success of reinforcement learning and outline several promising avenues for enhancing reinforcement learning in LLMs.
\end{itemize}

\subsection{Summary of Evaluations and Metrics}

\begin{itemize}[topsep=0pt]
    \item \textbf{English and Chinese Mathematical Reasoning}: Our models are systematically evaluated on a broad range of English and Chinese benchmarks, encompassing mathematical questions from elementary to collegiate complexity. For English, we utilize datasets such as GSM8K \citep{gsm8k}, MATH \citep{MATH}, SAT \citep{llemma}, OCW Courses \citep{minerva}, and MMLU-STEM \citep{mmlu}. On the Chinese side, we assess performance on MGSM-zh \citep{mgsm}, CMATH \citep{wei2023cmath}, Gaokao-MathCloze \citep{agieval}, and Gaokao-MathQA \citep{agieval}. The evaluation focuses on both the models' capacity to construct stand-alone text-based solutions and their proficiency in leveraging Python for problem-solving.

On standard English evaluation sets, DeepSeekMath-Base achieves results on par with the proprietary Minerva 540B \citep{minerva} model, while outperforming all publicly available base models—such as Mistral 7B \citep{mistral} and Llemma-34B \citep{llemma}—by a notable margin, irrespective of their prior math-specific pre-training. Significantly, DeepSeekMath-Base also excels on Chinese evaluation sets. This advantage is likely attributable to our approach of incorporating high-quality non-English mathematical data in pre-training, rather than restricting the data to English alone as previous studies \citep{minerva,llemma} have done. After applying mathematical instruction tuning and reinforcement learning methods, DeepSeekMath-Instruct and DeepSeekMath-RL attain robust results, achieving greater than 50\% accuracy on the challenging, competition-grade MATH dataset—a milestone reached for the first time among open-source systems.

\item \textbf{Formal Mathematics}: We assess DeepSeekMath-Base on the informal-to-formal theorem proving benchmark outlined by \citep{dsp_proof}, utilizing the miniF2F dataset \citep{minif2f} with the Isabelle proof assistant \citep{isabelle}. Our evaluation indicates that DeepSeekMath-Base achieves robust few-shot performance in autoformalization tasks.
\item \textbf{Natural Language Understanding, Reasoning, and Code}: To thoroughly measure the model's capabilities in general comprehension, logical reasoning, and programming, we subject DeepSeekMath-Base to the Massive Multitask Language Understanding (MMLU) benchmark \citep{mmlu}, which features 57 diverse multiple-choice tasks, and to BIG-Bench Hard (BBH) \citep{bbh}, comprising 23 rigorously challenging problems that often demand multi-step reasoning. Additionally, the model is evaluated on HumanEval \citep{codex} and MBPP \citep{mbpp}, both standard benchmarks for code generation. Pre-training with mathematics data notably enhances the model’s performance in both linguistic understanding and reasoning tasks.

\section{Math Pre-Training}

\subsection{Data Collection and Decontamination} This section details the methodology used to assemble the DeepSeekMath Corpus utilizing data from Common Crawl. Figure \ref{fig:data_collect} illustrates an iterative workflow designed to efficiently collect a substantial mathematical dataset, beginning with a seed corpus (such as a curated, high-quality set of math-related data). It should be emphasized that this strategy is equally transferable to additional areas, for example, programming.

Initially, we select OpenWebMath \citep{openwebmath}, which contains curated, high-quality mathematical web texts, to serve as our foundational seed dataset. Leveraging this dataset, we train a fastText model \citep{joulin2016fasttext} to retrieve additional web pages with characteristics similar to those in OpenWebMath. In particular, we randomly extract 500,000 instances from the seed corpus as positive examples, paired with a further 500,000 web pages from Common Crawl designated as negative samples. The training process utilizes an open-source implementation\footnote{\url{https://fasttext.cc}}, setting the vector size at 256, the learning rate at 0.1, the maximum n-gram size to 3, a minimum threshold of 3 word occurrences, and running for 3 epochs. To manage the volume of data from Common Crawl, we apply deduplication based on URLs and perform near-duplicate detection, which trims the dataset to 40 billion HTML web pages. Subsequently, the trained fastText model is employed to retrieve mathematical web documents from this deduplicated corpus. In order to eliminate content of inferior quality, we prioritize the identified documents by their fastText-assigned scores, retaining only the highest-scoring pages. The final selection volume is determined through preliminary pre-training on subsets of the top 40B, 80B, 120B, and 160B tokens. For this initial round, we opt to retain the subset comprising the top 40 billion tokens.

Following the initial round of data acquisition, a significant number of mathematical web pages remain outside the collected set. This is largely attributable to the limited heterogeneity in the positive instances used to train the fastText model. To address this, we identify new sources of mathematical web content and integrate them into the seed dataset, aiming to enhance the fastText model’s coverage. Our process begins by partitioning the full Common Crawl dataset into mutually exclusive domains, with each domain comprising web pages that share a common base URL. For each of these domains, we determine the proportion of web pages gathered during the first round of data collection. Any domain for which more than 10\% of the content has been retrieved is flagged as being math-related (for instance, \url{mathoverflow.net}).
Next, we carry out manual annotation of URLs within these math-focused domains that clearly correspond to mathematical material (e.g., \url{mathoverflow.net/questions}).
Uncollected web pages that correspond to these manually annotated URLs are then incorporated into the seed dataset. This methodology allows for an expansion of the positive example set, resulting in the training of a more robust fastText model, which can identify a greater volume of mathematical web pages in future iterations. Upon completing four cycles of data collection, we obtain a final dataset consisting of 35.5M mathematical web pages, which amounts to 120B tokens. By the fourth cycle, we observe that approximately 98\% of the mathematical data has already been retrieved by the previous iteration, prompting us to conclude the data collection process.

In order to prevent benchmark contamination, we adopt the approach from \cite{deepseek-coder} by excluding web pages that feature questions or answers from established English mathematical benchmarks, including GSM8K \citep{gsm8k} and MATH \citep{MATH}, as well as Chinese datasets such as CMATH \citep{wei2023cmath} and AGIEval \citep{agieval}. The procedure for filtering is as follows: any segment of text containing a 10-gram phrase that exactly matches a sub-string found in any of these benchmark datasets is eliminated from our math training data. Additionally, for cases where benchmark passages are shorter than 10 grams but comprise at least 3 grams, we use exact string matching to identify and exclude any web pages containing such content.

\subsection{Validating the Quality of the DeepSeekMath~Corpus}

We run pre-training experiments to investigate how the DeepSeekMath~Corpus is compared with the recently released math-training corpora:

\begin{itemize}[topsep=0pt]
    \item \textbf{MathPile} \citep{mathpile}: This dataset is composed of multiple sources, amassing 8.9B tokens from resources such as textbooks, Wikipedia, ProofWiki, CommonCrawl, StackExchange, and arXiv, with arXiv constituting the predominant share (exceeding 85\%).
    \item \textbf{OpenWebMath} \citep{openwebmath}: This corpus includes 13.6B tokens extracted from CommonCrawl, selected through filtration techniques targeting mathematical material.
    \item \textbf{Proof-Pile-2} \citep{llemma}: This dataset integrates OpenWebMath, AlgebraicStack (containing 10.3B tokens of code relevant to mathematics), and arXiv literature (with 28.0B tokens). For experiments involving Proof-Pile-2, we adhere to the ratio of arXiv:Web:Code as 2:4:1 as specified in \cite{llemma}.
\end{itemize}

\subsubsection{Training Setting}
\label{sec:quality-policy}
We conduct mathematical fine-tuning on a general-purpose language model containing 1.3B parameters, constructed using the same architecture as DeepSeek LLMs \citep{deepseek-llm}, and refer to this model as DeepSeek-LLM 1.3B. For each separate mathematical dataset, the model is trained over 150B tokens. The entirety of our experiments is implemented via the lightweight and efficient HAI-LLM framework \citep{haillm}. In alignment with the DeepSeek LLMs' training strategies, we employ the AdamW optimizer \citep{adamW} configured with $\beta_1=0.9$, $\beta_2=0.95$, and $\mathrm{weight\_decay}=0.1$. The learning rate is controlled by a multi-stage schedule: it rises to a maximum after 2,000 warmup steps, decays to 31.6\% of the maximum at 80\% of total training, and drops further to 10.0\% at the 90\% completion mark. The peak learning rate is set to 5.3e-4. We adopt a batch size comprising 4M tokens and utilize a context window of 4K tokens.

\subsubsection{Evaluation Results}

\textbf{The DeepSeekMath~Corpus is of high quality, covers multilingual mathematical content, and is the largest in size.}

\begin{itemize}[topsep=0pt]
    \item \textbf{High-quality}:
    To assess downstream capabilities, we measured few-shot chain-of-thought performance across 8 mathematical benchmarks \cite{cot}. Table \ref{tab:corpora_comparison} illustrates that models fine-tuned with the DeepSeekMath~Corpus outperform others. Furthermore, Figure \ref{fig:corpora_comparisons} reveals that, after processing 50B tokens (equivalent to a single pass over Proof-Pile-2), models trained on DeepSeekMath display superior results relative to those trained on Proof-Pile-2, suggesting that the overall quality of DeepSeekMath Corpus is notably higher.
    \item \textbf{Multilingual}:
    The DeepSeekMath~Corpus contains content in several languages, with English and Chinese constituting the most substantial portions. According to Table \ref{tab:corpora_comparison}, utilizing DeepSeekMath~Corpus as training material advances mathematical reasoning ability in both English and Chinese. This contrasts with prior mathematical corpora, which focus predominantly on English and therefore offer marginal or negative effects when applied to tasks requiring Chinese mathematical reasoning.
    \item \textbf{Large-scale}:
    In terms of size, the DeepSeekMath~Corpus far exceeds current mathematical corpora. Figure \ref{fig:corpora_comparisons} demonstrates that DeepSeek-LLM 1.3B, trained with DeepSeekMath~Corpus, experiences accelerated and sustained gains. Meanwhile, baseline datasets—being much more limited in scale—are recycled multiple times during training, which causes their model improvements to stagnate rapidly.
\end{itemize}

\subsection{Training and Evaluating DeepSeekMath-Base 7B}

In this section, we present DeepSeekMath-Base 7B, a foundational model engineered for robust reasoning, particularly within mathematical domains. The model’s initial parameters are taken from DeepSeek-Coder-Base-v1.5 7B \citep{deepseek-coder}, after which it undergoes training on a dataset of 500B tokens. The data is allocated as follows: 56\% originates from the DeepSeekMath Corpus, 4\% from AlgebraicStack, 10\% from arXiv, 20\% consists of Github code, and the final 10\% is comprised of natural language samples derived from Common Crawl in both English and Chinese. Training largely follows the procedures detailed in Section \ref{sec:quality-policy}; however, the maximum learning rate is adjusted to 4.2e-4 and the batch size is set at 10M tokens.

We present an extensive evaluation of DeepSeekMath-Base 7B’s mathematical proficiency, with particular emphasis on its competence in generating independent mathematical solutions unaided by external resources, employing tools to tackle mathematical problems, and performing formal theorem verification. Additionally, our analysis extends to a broader characterization of the base model’s abilities, covering aspects such as natural language comprehension, logical reasoning, and programming aptitude.

\paragraph{Mathematical Problem Solving with Step-by-Step Reasoning}

We assess the capability of DeepSeekMath-Base in tackling mathematical tasks through few-shot chain-of-thought prompting \citep{cot} on eight distinct English and Chinese benchmarks. These benchmarks span areas of quantitative reasoning—such as GSM8K \citep{gsm8k}, MATH \citep{MATH}, and CMATH \citep{wei2023cmath}—as well as multiple-choice assessments like MMLU-STEM \citep{mmlu} and Gaokao-MathQA \citep{agieval}. The evaluation covers a broad spectrum of mathematical disciplines, ranging from elementary to college-level problem sets.

According to the results summarized in Table \ref{tab:base_math_cot}, DeepSeekMath-Base 7B consistently achieves the top scores across all eight evaluated benchmarks among open-source base models—this group includes both the commonly adopted general-purpose model Mistral 7B \citep{mistral} and the newer Llemma 34B \citep{llemma}, which was specifically trained on mathematical content from Proof-Pile-2 \citep{llemma}. Particularly striking is DeepSeekMath-Base’s performance on the challenging, competition-oriented MATH dataset: it exceeds the next-best open-source base models by a margin greater than 10\% in absolute terms. Furthermore, it even outperforms Minerva 540B \citep{minerva}, a proprietary base model built on PaLM \citep{palm}, further refined with mathematical data, and with a parameter count 77 times larger.

\paragraph{Mathematical Problem Solving with Tool Use} To assess the mathematical reasoning abilities of models augmented with tool usage, we perform evaluations on GSM8K and MATH benchmarks employing a few-shot program-of-thought prompting approach \citep{pot,pal}. For each task, models are instructed to generate a Python script to address the problem, and are permitted to leverage libraries like \textit{math} and \textit{sympy} for advanced calculations. The answer is then derived from the output produced by executing these generated programs. According to Table \ref{tab:base_math_pot_proof}, DeepSeekMath-Base 7B demonstrates superior performance compared to the previous leading model, Llemma 34B.

\paragraph{Formal Mathematics}

The automation of formal proofs plays a pivotal role in guaranteeing both the correctness and robustness of mathematical arguments, while also streamlining the proving process—a trend that has garnered growing interest recently. In this study, we assess DeepSeekMath-Base 7B on the informal-to-formal proof generation task outlined in \citep{dsp_proof}, which involves producing a formalized proof given an informal description, its corresponding formal statement, and an informal reasoning sequence. For our experiments, we utilize the miniF2F benchmark \citep{minif2f}, which focuses on formalized, Olympiad-level mathematical challenges. Specifically, we apply few-shot prompting techniques to produce formal proofs in Isabelle for each benchmark problem. In alignment with the protocol established by \cite{dsp_proof}, our methodology employs the model to generate high-level proof outlines, subsequently completing the finer details using the Sledgehammer automated theorem prover \citep{sledgehammer}. According to the results displayed in Table \ref{tab:base_math_pot_proof}, DeepSeekMath-Base 7B exhibits considerable effectiveness in the task of proof autoformalization.

\paragraph{Natural Language Understanding, Reasoning, and Code}

We assess the performance of our model in natural language understanding using MMLU \citep{mmlu}, in reasoning tasks via BBH \citep{bbh}, and in coding proficiency on the HumanEval \citep{codex} and MBPP \citep{mbpp} benchmarks. Table \ref{tab:understanding_reasoning_code} demonstrates that DeepSeekMath-Base 7B achieves marked improvements on MMLU and BBH in comparison to its predecessor, DeepSeek-Coder-Base-v1.5 \citep{deepseek-coder}, highlighting the beneficial effects of mathematical training for language understanding and reasoning tasks. Moreover, through continued training that incorporates code tokens, DeepSeekMath-Base 7B preserves the strong coding performance seen in DeepSeek-Coder-Base-v1.5 across both coding benchmarks. In summary, DeepSeekMath-Base 7B provides a substantial performance advantage over the generalist Mistral 7B model \citep{mistral} across all three evaluated reasoning and coding tasks.

\section{Supervised Fine-Tuning}

\subsection{SFT Data Curation}

We develop a mathematically focused instruction-tuning dataset that spans both English and Chinese questions, encompassing a range of mathematical domains and difficulty tiers. Each problem is annotated with a corresponding solution expressed in one of several formats: chain-of-thought (CoT) \citep{cot}, program-of-thought (PoT) \citep{pot,pal}, or tool-integrated reasoning \citep{tora}. The dataset comprises a total of 776,000 training examples.

\begin{itemize}[topsep=0pt]
    \item \textbf{English mathematical datasets}: We provide tool-supported solution annotations for both GSM8K and MATH datasets. In addition, we utilize a selection from MathInstruct \citep{MathInstruct} and include the training split of Lila-OOD \citep{lila}, all featuring solutions formulated using either CoT or PoT methods. Our assembled English datasets encompass a broad spectrum of mathematical disciplines, including algebra, probability, number theory, calculus, and geometry.
    \item \textbf{Chinese mathematical datasets}: Our dataset comprises Chinese K-12 math questions that cover 76 different sub-topics (such as linear equations), each paired with explanations generated through both CoT reasoning and tool-integrated approaches.
\end{itemize}

\subsection{Training and Evaluating DeepSeekMath-Instruct 7B}

This section presents DeepSeekMath-Instruct 7B, a model fine-tuned for mathematical tasks using instruction tuning, with DeepSeekMath-Base serving as its foundation. During training, instances are concatenated in a random order, up to a total context size limit of 4K tokens. The model undergoes training for 500 iterations, utilizing a batch size of 256 and maintaining a fixed learning rate of 5e-5 throughout the process.

We evaluate models' mathematical performance both without and with tool use, on 4 quantitative reasoning benchmarks in English and Chinese. We benchmark our model against the leading models of the time:

\begin{itemize}[topsep=0pt]
    \item \textbf{Closed-source models} comprise: (1) the GPT series, with GPT-4 \citep{gpt4} and its variant, GPT-4 Code Interpreter~\footnote{\url{https://openai.com/blog/chatgpt-plugins\#\#code-interpreter}}, recognized as the most advanced; (2) Gemini Ultra and Gemini Pro \citep{gemini}; (3) Inflection-2 \citep{inflection-2}; (4) Grok-1~\footnote{\url{https://x.ai/model-card}}; in addition, models launched by Chinese technology firms such as (5) Baichuan-3~\footnote{\url{https://www.baichuan-ai.com}} and (6) the new GLM-4~\footnote{\url{https://open.bigmodel.cn/dev/api\#glm-4}}

    from the GLM family \citep{glm}.

These models are designed for broad applications, and the majority have been subjected to various alignment processes.
    \item \textbf{Open-source models} encompass both general-purpose and mathematics-enhanced variants: examples of general models include (1) DeepSeek-LLM-Chat 67B \citep{deepseek-llm}, (2) Qwen 72B \citep{qwen}, (3) SeaLLM-v2 7B \citep{seallm}, and (4) ChatGLM3 6B \citep{chatglm3}. Mathematical reasoning models comprise (5) InternLM2-Math 20B~\footnote{\url{https://github.com/InternLM/InternLM-Math}}, an extension of InternLM2 that has received specialized mathematics instruction and subsequent instruction tuning, (6) Math-Shepherd-Mistral 7B, which applies PPO training \citep{schulman2017proximal} and employs a process-supervised reward model to Mistral 7B \citep{mistral}, (7) the WizardMath series \citep{wizardmath}, which enhances Mistral 7B and Llama-2 70B \citep{llama2} for mathematical tasks through evolve-instruct (an approach to instruction tuning utilizing AI-evolved prompts) and PPO, leveraging training data mainly from GSM8K and MATH, (8) MetaMath 70B \citep{metamath}, a Llama-2 70B version refined on expanded GSM8K and MATH datasets, (9) ToRA 34B \cite{tora}, which represents a fine-tuned version of CodeLlama 34B focused on tool-based mathematical reasoning, and (10) MAmmoTH 70B \citep{MathInstruct}, which fine-tunes Llama-2 70B with MathInstruct via instruction tuning.

Table \ref{tab:sft_rl_math} presents results for the evaluation setting without tool assistance, highlighting the strong step-by-step reasoning abilities of DeepSeekMath-Instruct 7B. Remarkably, our model exceeds the performance of all open-source models as well as most commercial systems (such as Inflection-2 and Gemini Pro) on the challenging, competition-level MATH dataset, achieving at least a 9\% absolute margin. This outperformance persists even in comparison with models possessing significantly more parameters (e.g., Qwen 72B) or those optimized using math-centric reinforcement learning techniques (such as WizardMath-v1.1 7B). While DeepSeekMath-Instruct reaches parity with Chinese commercial offerings like GLM-4 and Baichuan-3 in MATH performance, its results remain lower than those of GPT-4 and Gemini Ultra.

When assessed in a scenario permitting the use of both natural language reasoning and programmatic tools to solve problems, DeepSeekMath-Instruct 7B attains nearly 60\% accuracy on the MATH benchmark, outperforming all previously released open-source models. Across additional evaluation datasets, our model's performance rivals that of DeepSeek-LLM-Chat 67B—the previous leading model that possesses ten times more parameters.  
\section{Reinforcement Learning}

\subsection{Group Relative Policy Optimization} Reinforcement learning (RL) has demonstrated its capacity to enhance the mathematical reasoning skills of LLMs subsequent to the Supervised Fine-Tuning (SFT) phase \citep{wang2023math,wizardmath}. In the following, we present Group Relative Policy Optimization (GRPO), a novel RL algorithm that offers both efficiency and efficacy.

\subsubsection{From PPO to GRPO} Proximal Policy Optimization (PPO) \citep{schulman2017proximal} is an actor-critic RL algorithm that is widely used in the RL fine-tuning stage of LLMs \citep{ouyang2022training}. In particular, it optimizes LLMs by maximizing the following surrogate objective:

\begin{equation}
\footnotesize
    \mathcal{J}_{PPO}(\theta) = \mathbb{E}{[q \sim P(Q), o \sim \pi_{\theta_{old}}(O|q)]} \frac{1}{|o|} \sum_{t=1}^{|o|} \min \left[ \frac{\pi_\theta(o_{t} | q, o_{<t})}{\pi_{\theta_{old}}(o_{t} | q, o_{<t})} A_{t}, \text{clip} \left( \frac{\pi_\theta(o_{t} | q, o_{<t})}{\pi_{\theta_{old}}(o_{t} | q, o_{<t})}, 1 - \varepsilon, 1 + \varepsilon \right)  A_{t} \right] ,
\end{equation}

The equation above describes the loss objective for Proximal Policy Optimization (PPO), in which the expectation is computed over queries sampled from $P(Q)$ and corresponding outputs drawn from the policy $\pi_{\theta_{old}}(O|q)$. At each position $t$ within the sequence $o$, the loss includes an average over the sequence length and takes the minimum between the policy ratio multiplied by the advantage $A_t$, and a clipped version of this ratio bounded between $1 - \varepsilon$ and $1 + \varepsilon$, also scaled by $A_t$. This mechanism serves to restrict large, destabilizing policy updates during optimization.

Here, $\pi_{\theta}$ and $\pi_{\theta_{old}}$ denote the present and previous policy models, while $q$ and $o$ correspond to questions drawn from the dataset and responses generated by the earlier policy $\pi_{\theta_{old}}$. The term $\varepsilon$ refers to the clipping hyperparameter used in PPO to enhance the robustness of training. The advantage term $A_t$ is estimated using Generalized Advantage Estimation (GAE) \citep{gae}, which utilizes the future rewards $\{r_{\ge t}\}$ and an approximated value function $V_{\psi}$. Consequently, in PPO, both the policy and value functions are jointly learned. To prevent excessive optimization towards the reward model, it is common to incorporate a KL divergence penalty at every token in the reward, calculated with respect to a reference model, as suggested by \citep{ouyang2022training}, specifically,

\begin{equation}
    r_{t} = r_\varphi(q, o_{\le t}) - \beta \log\frac{\pi_{\theta}(o_{t}|q, o_{<t})}{\pi_{ref}(o_{t}|q, o_{<t})},
\label{eq:PPO-reward}
\end{equation}

Here, $r_\varphi$ denotes the reward model, while $\pi_{ref}$ refers to the reference model—typically the original SFT model—with $\beta$ serving as the parameter that weights the KL divergence penalty.

As the value function employed in PPO is typically another model of comparable size as the policy model, it brings a substantial memory and computational burden. Additionally, during RL training, the value function is treated as a baseline in the calculation of the advantage for variance reduction. While in the LLM context, usually only the last token is assigned a reward score by the reward model, which may complicate the training of a value function that is accurate at each token. To address this, as shown in Figure \ref{fig:grpo}, we propose Group Relative Policy Optimization (GRPO), which obviates the need for additional value function approximation as in PPO, and instead uses the average reward of multiple sampled outputs, produced in response to the same question, as the baseline. More specifically, for each question $q$, GRPO samples a group of outputs $\{o_1, o_2, \cdots, o_G\}$  from the old policy  $\pi_{\theta_{old}}$  and then optimizes the policy model by maximizing the following objective:

\begin{equation}
\footnotesize
\begin{split}
    \mathcal{J}_{GRPO}(\theta) &= \mathbb{E}_{q \sim P(Q), \{o_i\}_{i=1}^G \sim \pi_{\theta_{old}}(O|q)}  \\
    & \frac{1}{G}\sum_{i=1}^G\frac{1}{|o_i|} \sum_{t=1}^{|o_i|} \Bigg\{ \min \left[ \frac{\pi_\theta(o_{i,t} | q, o_{i,<t})}{\pi_{\theta_{old}}(o_{i,t} | q, o_{i,<t})} \hat{A}_{i,t},\ \text{clip}\left( \frac{\pi_\theta(o_{i,t} | q, o_{i,<t})}{\pi_{\theta_{old}}(o_{i,t} | q, o_{i,<t})},\ 1 - \varepsilon,\ 1 + \varepsilon \right) \hat{A}_{i,t} \right] \\
    & \hspace{6.4cm} - \beta\ \mathbb{D}_{KL} \left[ \pi_{\theta} \,\|\, \pi_{ref} \right] \Bigg\},
\end{split}
\label{eq:GRPO-obj}
\end{equation}

where $\varepsilon$ and $\beta$ are hyper-parameters, and $\hat{A}_{i,t}$  is the advantage calculated based on relative rewards of the outputs inside each group only, which will be detailed in the following subsections. The group relative way that GRPO leverages to calculate the advantages, aligns well with the comparative nature of rewards models,  as reward models are typically trained on datasets of comparisons between outputs on the same question. Also note that, instead of adding KL penalty in the reward, GRPO regularizes by directly adding the KL divergence between the trained policy and the reference policy to the loss, avoiding complicating the calculation of  $\hat{A}_{i,t}$. And different from the KL penalty term used in (\ref{eq:PPO-reward}), we estimate the KL divergence with the following unbiased estimator \citep{kl_approx}:

\begin{equation}
\small
    \mathbb{D}_{KL}\left[\pi_{\theta} || \pi_{ref}\right] = \frac{\pi_{ref}(o_{i,t}|q,o_{i,<t})}{\pi_{\theta}(o_{i,t}|q,o_{i,<t})}- \log\frac{\pi_{ref}(o_{i,t}|q,o_{i,<t})}{\pi_{\theta}(o_{i,t}|q,o_{i,<t})} - 1,
\end{equation}

which is guaranteed to be positive.

\begin{algorithm}[t]
  \small
  \caption{Iterative Group Relative Policy Optimization}
  \textbf{Input} initial policy model $\pi_{\theta_{\text{init}}}$; reward models $r_\varphi$; task prompts $\mathcal{D}$; 
  hyperparameters $\varepsilon$, $\beta$, $\mu$
  \begin{algorithmic}[1]
    \State Initialize policy model: $\pi_\theta \leftarrow \pi_{\theta_{\text{init}}}$
    \For{iteration = 1, \dots, I}
       \State Assign reference model $\pi_{ref} \leftarrow \pi_{\theta}$
      \For{step = 1, \dots, M}
      \State Draw a mini-batch $\mathcal{D}_b$ from $\mathcal{D}$
      \State Save current policy as $\pi_{\theta_{old}} \leftarrow \pi_{\theta}$ 
      \State For every query $q \in \mathcal{D}_b$, generate $G$ responses $\{o_i\}_{i=1}^G \sim \pi_{\theta_{old}} (\cdot \mid q)$
      \State Assign corresponding rewards $\{r_i\}_{i=1}^{G}$ to each sampled output $o_i$ with $r_{\varphi}$
      \State Apply group relative advantage estimation to calculate $\hat{A}_{i,t}$ for each token $t$ in every $o_i$
      \For{GRPO iteration = 1, \dots, $\mu$}
        \State Adjust policy parameters $\pi_{\theta}$ by maximizing the GRPO objective (Equation \ref{eq:GC-GRPO})
      \EndFor
    \EndFor \State Perform ongoing reward model training on $r_\varphi$ by employing a replay buffer. 
    \EndFor 
  \end{algorithmic}
  \textbf{Output} $\pi_\theta$
  \label{alg:iter-grpo}
\end{algorithm}

\subsubsection{Outcome Supervision RL with GRPO} In this approach, for any given question $q$, a collection of outputs $\{o_1, o_2, \cdots, o_G\}$ is generated using the previous iteration of the policy model, denoted as $\pi_{\theta_{old}}$. Each of these outputs is evaluated with a reward model, resulting in $G$ corresponding rewards $\mathbf{r} = \{r_1, r_2, \cdots, r_G\}$. The obtained reward values are subsequently standardized by removing the mean and scaling by the standard deviation within the sampled group. In the context of outcome supervision, each output $o_i$ is assigned this normalized terminal reward, and all token-level advantages $\hat{A}_{i, t}$ throughout $o_i$ are set equal to the same normalized reward, that is, $\hat{A}_{i, t} = \widetilde{r}_i = \frac{r_i- {\rm mean}(\mathbf{r})}{{\rm std}(\mathbf{r})}$. The policy is then updated to optimize the objective given in equation (\ref{eq:GRPO-obj}).

\subsubsection{Process Supervision RL with GRPO} Outcome supervision delivers feedback only at the conclusion of each generated output, which may be inadequate for effectively guiding the policy in intricate mathematical problem settings. In alignment with \cite{wang2023math}, we further investigate process supervision, a strategy that allocates reward at each stage of intermediate reasoning. Specifically, let the question $q$ be given and suppose $G$ output sequences $\{o_1, o_2, \cdots, o_G\}$ are sampled. A process reward model evaluates each reasoning step across all outputs, resulting in reward sets: $\mathbf{R} = \{ \{r_1^{index(1)},\cdots,r_1^{index(K_1)}\}, \cdots,  \{r_G^{index(1)},\cdots,r_G^{index(K_G)}\} \}$, with $index(j)$ denoting the position of the terminating token of the $j$-th step and $K_i$ indicating the number of steps for output $i$. To standardize these evaluations, the rewards are normalized using the mean and standard deviation, as $\widetilde{r}_i^{index(j)} = \frac{r_i^{index(j)} - {\rm mean(\mathbf{R})}}{{\rm std(\mathbf{R})}}$.

The process supervision approach then computes the advantage for every token as the cumulative sum of normalized rewards for all succeeding steps, i.e., $\hat{A}_{i, t} = \sum_{index(j) \ge t} \widetilde{r}_i^{index(j)}$.
Finally, the optimization of the policy is carried out by maximizing the objective stated in equation (\ref{eq:GRPO-obj}).

\subsubsection{Iterative RL with GRPO} During the course of reinforcement learning, reliance on an outdated reward model may not provide adequate supervision for the evolving policy model. To address this, we investigate iterative RL utilizing GRPO. As detailed in Algorithm \ref{alg:iter-grpo}, iterative GRPO involves assembling new reward model training datasets by sampling from the policy model and updating the previous reward model via a replay strategy that retains 10\% of prior samples. Subsequently, the policy model is established as the reference model, and it is further optimized using the updated reward model.

\subsection{Training and Evaluating DeepSeekMath-RL}

Our reinforcement learning (RL) experiments are carried out with DeepSeekMath-Instruct 7B as the backbone. For the RL training stage, we utilize questions formatted with chain-of-thought reasoning, sourced from the GSM8K and MATH subsets within our supervised fine-tuning (SFT) dataset, yielding a corpus of approximately 144K items. We intentionally omit SFT questions outside these domains to specifically analyze how RL influences benchmark results when relevant data is absent during RL training. The assembly of reward model training data mirrors the methodology described in \citep{wang2023math}. The initial reward model is initialized from DeepSeekMath-Base 7B and optimized using a learning rate of 2e-5. For the GRPO algorithm, the policy model’s learning rate is set at 1e-6, and a KL coefficient of 0.04 is applied. Each question is used to generate $64$ samples. The input sequence length is capped at 1024 tokens, and we set the batch size to 1024 for training. After each exploration phase, only a single policy update is performed. We follow the DeepSeekMath-Instruct 7B evaluation setup to assess DeepSeekMath-RL 7B’s performance across benchmarks. Within the context of DeepSeekMath-RL 7B, GSM8K and MATH questions requiring chain-of-thought solutions are categorized as in-domain, while all remaining benchmarks are treated as out-of-domain.

Table \ref{tab:sft_rl_math} presents comparative results for both open- and closed-source models on English and Chinese benchmarks, employing chain-of-thought and tool-augmented reasoning approaches. The observations are as follows: 1) Utilizing chain-of-thought reasoning, DeepSeekMath-RL 7B achieves accuracy rates of 88.2\% on GSM8K and 51.7\% on MATH. These scores not only exceed those of all open-source models within the 7B–70B parameter range but also outperform most closed-source systems. 2) Notably, DeepSeekMath-RL 7B’s training is limited strictly to chain-of-thought formatted instruction data for GSM8K and MATH, with its initialization derived from DeepSeekMath-Instruct 7B. In spite of this restricted training regime, it demonstrates superior performance relative to DeepSeekMath-Instruct 7B in all metrics evaluated, highlighting the advantages brought by reinforcement learning.

\section{Discussion}
Here, we present the results obtained from our pre-training and RL experimental investigations.

\subsection{Lessons Learnt in Pre-Training}

We begin by detailing our observations during the pre-training phase. Except where explicitly mentioned, all experiments follow the training configurations described in Section \ref{sec:quality-policy}. Specifically, any references to the DeepSeekMath Corpus in this section pertain to a dataset comprising 89 billion tokens, sourced from the second round of data collection.

\subsubsection{Code Training Benefits Mathematical Reasoning}

An often-cited but unsubstantiated claim posits that exposure to code enhances reasoning skills. In this work, we provide evidence addressing this claim in part, focusing specifically on mathematics: training on code bolsters models' performance in mathematical reasoning tasks, regardless of whether external tools are employed.

To study how code training affects mathematical reasoning, we experimented with the following two-stage training and one-stage training settings:

\textbf{Two-Stage Training}

\begin{itemize}[topsep=0pt]
    \item \textbf{Code Training for 400B Tokens $\rightarrow$ Math Training for 150B Tokens}: The DeepSeek-LLM 1.3B model is first trained on a dataset consisting of 400B code tokens, after which it is further trained on 150B math tokens.
    \item \textbf{General Training for 400B Tokens $\rightarrow$ Math Training for 150B Tokens}: For comparison, we conduct a parallel experiment where the initial stage uses 400B general tokens—drawn from a broad general-purpose corpus assembled by DeepSeek-AI—rather than code tokens, to assess whether code data provides benefits over general language data in mathematical reasoning capabilities.
\end{itemize}

\textbf{One-Stage Training}

\begin{itemize}[topsep=0pt]
    \item \textbf{Math Training over 150B Tokens}: The DeepSeek-LLM 1.3B model undergoes training exclusively on 150B tokens comprised of mathematical content;
    \item \textbf{Joint Training with 400B Code Tokens and 150B Math Tokens}: Post-code training math fine-tuning has been observed to reduce code task performance. Thus, we explore a single-stage training approach where code and math tokens are jointly used, aiming to both enhance mathematical reasoning capabilities and mitigate catastrophic forgetting.
\end{itemize}

\paragraph{Results}
The downstream outcomes corresponding to various training configurations are presented in Table \ref{tab:code-math} and Table \ref{tab:code-math-general-eval}.

Training models with code yields substantial improvements in program-aided mathematical reasoning, applicable to both two-stage and one-stage training paradigms. As evidenced in Table \ref{tab:code-math}, incorporating code training during the initial stage of the two-stage setup markedly boosts performance on GSM8K and MATH benchmarks in Python. Introducing dedicated math training in the subsequent stage further elevates these results. Notably, when employing a one-stage training approach, simultaneously integrating code and math tokens not only alleviates the catastrophic forgetting commonly observed in two-stage procedures but also leads to complementary gains in both code generation (Table \ref{tab:code-math-general-eval}) and program-enhanced mathematical reasoning (Table \ref{tab:code-math}).

Training with code also enhances mathematical reasoning abilities even in the absence of tool use. Within the two-stage training framework, the preliminary phase of code training yields moderate gains on its own. Furthermore, it increases the effectiveness of the following mathematical training, which ultimately achieves optimal results. In contrast, merging code tokens and mathematical tokens for single-stage training hinders mathematical reasoning in contexts without tool use. A plausible explanation is that DeepSeek-LLM 1.3B, given its relatively small model size, does not possess sufficient capacity to effectively integrate both code and mathematical information at the same time.

\subsubsection{ArXiv Papers Seem Ineffective in Improving Mathematical Reasoning}

ArXiv papers are commonly included as a component of math pre-training data \citep{minerva,gpt-f,llemma,mathpile}. However, detailed analysis regarding their impact on mathematical reasoning has not been extensively conducted. Perhaps counter-intuitively, according to our experiments, arXiv papers seem ineffective in improving mathematical reasoning. We experiment with models of different sizes, including DeepSeek-LLM 1.3B and DeepSeek-Coder-Base-v1.5 7B \citep{deepseek-coder}, using arXiv corpora that underwent varied processing pipelines:

\begin{itemize}[topsep=0pt]
    \item \textbf{MathPile} \citep{mathpile}: a dataset containing 8.9 billion tokens, primarily composed (over 85\%) of scientific papers from arXiv, assembled through the application of heuristic approaches for data cleaning and filtering;
    \item \textbf{ArXiv-RedPajama} \citep{redpajama}: a corpus comprising all available arXiv LaTeX documents with extraneous elements such as preambles, comments, macros, and bibliographies systematically stripped out, resulting in a total size of 28.0 billion tokens.
\end{itemize}

In our study, we individually pre-train DeepSeek-LLM 1.3B with 150B tokens and DeepSeek-Coder-Base-v1.5 7B with 40B tokens using distinct arXiv corpora. Our findings suggest that arXiv papers contribute little to enhancing mathematical reasoning capabilities. When models are exposed exclusively to arXiv datasets, neither demonstrates noticeable progress—in fact, performance sometimes regresses—on a range of mathematical benchmarks utilized here. These benchmarks span quantitative reasoning tasks such as GSM8K and MATH (Table \ref{tab:arxiv-cot}), multiple-choice assessments like MMLU-STEM (Table \ref{tab:arxiv-cot}), as well as formal mathematics tasks including miniF2F (Table \ref{tab:arxiv_atp}).

However, this conclusion has its limitations and should be taken with a grain of salt. We have not yet studied:

\begin{itemize}[topsep=0pt]
    \item How arXiv tokens influence particular mathematics-oriented applications not explored in this study, for instance, the informalization of theorems, which involves translating formal statements or proofs into their informal forms;
    \item The influence of arXiv tokens when integrated with alternative data sources;
    \item If scaling up model size would continue to highlight the advantages offered by arXiv papers.
\end{itemize}

Thus, further exploration is required, which we leave for future studies.

\subsection{Insights of Reinforcement Learning}
\subsubsection{Towards to a Unified Paradigm} Here, we introduce a cohesive framework to examine and compare various training approaches, including SFT, RFT, DPO, PPO, and GRPO. We also conduct empirical studies to investigate key components within this unified paradigm. In general, the gradient of a training method with respect to its parameter $\theta$ can be expressed as follows:

\begin{equation}
    \nabla_{\theta}\mathcal{J}_{\textcolor{red}{\mathcal{A}}}(\theta) = \mathbb{E}[\underbrace{(q,o) \sim \textcolor{red}{\mathcal{D}}}_{Data \ Source}]\left( \frac{1}{|o|} \sum_{t=1}^{|o|}  \underbrace{GC_{{\mathcal{A}}}(q, o, t, \textcolor{red}{\pi_{{rf}}})}_{Gradient \ Coefficient}  \nabla_{\theta}\log \pi_{\theta}(o_t | q, o_{<t})\right).
\label{eq:objective}
\end{equation}

There exist three key components: 1) \textit{Data Source $\mathcal{D}$}, which determines the training data; 2) \textit{Reward Function $\pi_{{rf}}$}, which is the source of the training reward signal; 3) \textit{Algorithm $\mathcal{A}$}: which processes the training data and the reward signal to the gradient coefficient $GC$ that determines the magnitude of the penalty or reinforcement for the data. We analyze several representative methods based on such a unified paradigm:

\begin{itemize}
    \item \textbf{Supervised Fine-tuning (SFT)}: SFT adapts a pretrained model by further training it on a curated set of human-annotated SFT data.
    \item \textbf{Rejection Sampling Fine-tuning (RFT)}: RFT builds upon the SFT model by fine-tuning it with selected outputs generated in response to SFT questions; only those outputs deemed correct are retained for training.
    \item \textbf{Direct Preference Optimization (DPO)}: DPO enhances the SFT model through additional fine-tuning, employing a pairwise DPO loss on extended outputs sampled from the SFT model itself.
    \item \textbf{Online Rejection Sampling Fine-tuning (Online RFT)}: Unlike standard RFT, Online RFT starts from the SFT model and continually refines the policy model using outputs sampled on-the-fly from the evolving policy model during training.
    \item \textbf{PPO/GRPO}: PPO/GRPO uses the SFT model as an initial policy and applies reinforcement by training on samples directly generated from the current policy model in real time.
\end{itemize}

Table \ref{tab:data_policy} provides an overview of the elements involved in these methods. For a comprehensive explanation of the derivation procedure, see Appendix \ref{app:analysis-rl}.

\paragraph{Observation about Data Source} We classify data sources into two main types: online sampling and offline sampling. In online sampling, training data is generated based on outputs from the current training policy model during ongoing exploration, whereas in offline sampling, training data is obtained from outputs of the pre-trained SFT model. Methods such as RFT and DPO utilize offline sampling, whereas Online RFT and GRPO are based on online sampling approaches.

As depicted in Figure \ref{fig:rl-analysis}, our results indicate that Online RFT achieves notably better performance than RFT across both benchmarks. In the initial phases of training, Online RFT performs on par with RFT; however, as training progresses, it obtains a pronounced performance lead, which underscores the benefits of an online learning approach. This observation aligns with expectations: during early training, the actor and SFT model remain quite similar, resulting in sampled data that only shows minimal distinctions. Conversely, as training advances, discrepancies in the data produced by the actor grow more prominent, making real-time data collection increasingly advantageous.

\paragraph{Observation about Gradient Coefficient} The algorithm utilizes the reward signal in determining the gradient coefficient, which is then used to update the model parameters. In our experiments, we categorize the reward function into `Rule' and `Model'. The `Rule' approach assesses response quality strictly by verifying answer correctness, whereas `Model' involves training a reward model that assigns a score to each response. Notably, the training data for the reward model are derived from the rule-based evaluations. As shown in Equations \ref{eq:GC-RFT} and \ref{eq:GC-GRPO}, there is a fundamental distinction between GRPO and Online RFT: GRPO modifies its gradient coefficient dynamically according to the reward value from the reward model, enabling it to differentially encourage or penalize responses based on the magnitude of their rewards. By comparison, Online RFT does not incorporate this adaptability—it treats all correct responses with an identical level of reinforcement and does not decrease reward for incorrect answers.

As illustrated in Figure \ref{fig:rl-analysis}, GRPO outperforms online RFT, underscoring the effectiveness of modifying the coefficients for positive and negative gradients. Moreover, GRPO+PS achieves better results than GRPO+OS, demonstrating the advantages of employing step-aware and fine-grained gradient coefficients. Additionally, iterative RL is investigated in our study, with experiments involving two successive iterations. Figure \ref{fig:iter-rl} reveals that iterative RL leads to substantial performance gains, most notably after the initial iteration.

\subsubsection{Why RL Works?} In this study, we apply reinforcement learning using a selected subset of instruction tuning data, and observe notable improvements in performance compared to the instruction-tuned baseline. To shed light on the underlying mechanisms for these gains, we compare the Pass@K and Maj@K accuracies of both Instruct and RL models across two benchmark datasets. According to Figure \ref{fig:maj-pass}, reinforcement learning leads to increased Maj@K scores but does not significantly impact Pass@K. This suggests that RL primarily strengthens the output distribution’s robustness, implying \textbf{the observed performance gains originate from promoting correct answers within the TopK responses rather than from advancing basic reasoning skills.} In a related vein, \citep{wang2023making} also discuss a \textbf{misalignment problem} encountered in reasoning tasks for SFT models, and demonstrate that applying various preference alignment approaches \citep{yuan2023rrhf,song2023preference,wang2023making} can enhance the reasoning capabilities of these models.

\subsubsection{How to Achieve More Effective RL?}
\label{sec:effec_rl}
Our findings indicate that RL is highly effective when applied to mathematical reasoning challenges. Moreover, we offer a coherent framework that enables a comprehensive understanding of prominent training strategies. Under this framework, each approach is viewed as either a form of direct or simplified RL. As outlined in Equation \ref{eq:objective}, three critical elements are involved: Data Source, Algorithm, and Reward Function. We also discuss several promising avenues for further advancement with respect to these components.

\paragraph{Data Source} The choice of data source underpins all forms of model training. Within the framework of RL, the term data source specifically denotes the set of unlabeled questions paired with responses generated by sampling from the policy model. In this study, our dataset consists exclusively of questions gathered during the instruction tuning phase, and outputs are generated through a straightforward nucleus sampling procedure. We posit that this limited approach may partially explain why our RL pipeline yields gains only in the Maj@K metric. Looking ahead, we aim to extend the evaluation of our RL pipeline to incorporate question prompts from outside the training distribution, alongside the use of \textbf{advanced sampling (decoding) strategies}, such as those employing tree-search algorithms \citep{yao2023tree}. Furthermore, we note that \textbf{efficient inference techniques} \citep{xia-etal-2023-speculative,leviathan2023fast,kwon2023efficient,xia2024unlocking}, which govern the exploration capabilities of policy models, constitute another critical factor in this context.

\paragraph{Algorithms} The algorithms interpret both the data and the reward feedback, leveraging this information through gradient coefficients in order to modify the model parameters. According to Equation \ref{eq:objective}, existing approaches generally place complete \textbf{TRUST} in the feedback from the reward function, using it as a decisive factor to modulate the conditional likelihood of specific tokens. Nonetheless, the dependability of this reward signal cannot always be guaranteed, particularly when confronted with highly complex tasks. As an illustration, the PRM800K datasets \citep{lightman2023let}—despite being meticulously curated by trained annotators—still exhibit nearly 20\% erroneous annotations\footnote{\url{https://github.com/openai/prm800k/issues/12\#issuecomment-1728491852}}. In response to this limitation, our focus will shift toward reinforcement learning algorithms designed to withstand noisy or unreliable reward feedback. We contend that the adoption of \textbf{WEAK-TO-STRONG} alignment techniques \citep{burns2023weak} could fundamentally reshape prevailing learning algorithm paradigms.

\paragraph{Reward Function} The reward function serves as the fundamental driver of the learning process. In RL contexts, it is typically parameterized by a neural reward model. We identify three pivotal avenues for advancing reward models: 1) \textbf{Strategies for improving the generalization capacity of reward models.} It is critical for reward models to robustly generalize to novel, out-of-distribution queries and sophisticated generation outputs; lacking this, reinforcement learning might only maintain the status quo of LLMs' output distributions without promoting deeper competency gains; 2) \textbf{Techniques for representing uncertainty within the reward model.} Capturing and utilizing model uncertainty can potentially create important synergies between imperfect reward modeling and algorithms designed to learn from weak supervision; 3) \textbf{Approaches for constructing high-fidelity process reward models efficiently}, enabling precise, granular feedback on the intermediate steps within the reasoning pipeline \citep{lightman2023let,wang2023math}.

\section{Conclusion, Limitation, and Future Work} In this work, we introduce DeepSeekMath, which surpasses all other open-source models on the competition-level MATH benchmark and achieves results close to those of proprietary systems. DeepSeekMath~is built upon DeepSeek-Coder-v1.5 7B as its foundation, followed by a continual pretraining phase using 500B tokens, incorporating 120B math-specific tokens predominantly extracted from Common Crawl. Through comprehensive ablation experiments, we demonstrate that web-based content constitutes a rich source of high-quality mathematical data, whereas data derived from arXiv does not yield the anticipated benefits. We propose Group Relative Policy Optimization (GRPO), an adaptation of Proximal Policy Optimization (PPO), that enhances mathematical reasoning ability while requiring less memory. Our empirical analysis indicates that GRPO can significantly improve performance, even when DeepSeekMath-Instruct 7B already achieves high benchmark scores. Finally, we present a cohesive framework to interpret various methodologies and highlight several promising research avenues to advance reinforcement learning techniques further.

While DeepSeekMath~demonstrates strong results on benchmarks assessing quantitative reasoning, its performance in areas such as geometry and theorem-proving remains lower compared to proprietary models. For example, during preliminary evaluations, DeepSeekMath~was unable to solve tasks involving triangles and ellipses, suggesting a possible bias in the types of data chosen during both pre-training and fine-tuning stages. Moreover, due to its model size constraints, DeepSeekMath~lags behind GPT-4 in few-shot learning scenarios; whereas GPT-4 shows marked improvement with few-shot examples, DeepSeekMath~exhibits nearly equivalent results under both zero-shot and few-shot settings. Moving forward, we intend to enhance our data engineering strategies to curate a higher-quality pre-training corpus. Additionally, we plan to investigate more effective approaches for LLM reinforcement learning, as outlined in Section \ref{sec:effec_rl}.

\end{document}